\documentclass[11pt]{article}
\usepackage[utf8]{inputenc}
\usepackage{amsmath,amssymb,amsthm}
\usepackage[margin=1in]{geometry}
\usepackage{hyperref}
\usepackage{xcolor}
\usepackage{enumitem}
\usepackage{newunicodechar}
\newunicodechar{ℝ}{\ensuremath{\mathbb{R}}}
\newunicodechar{ℕ}{\ensuremath{\mathbb{N}}}
\newunicodechar{⟨}{\ensuremath{\langle}}
\newunicodechar{⟩}{\ensuremath{\rangle}}
\newunicodechar{𝓝}{\ensuremath{\mathcal{N}}}
\newunicodechar{∂}{\ensuremath{\partial}}
\newunicodechar{·}{\ensuremath{\cdot}}
\newunicodechar{α}{\ensuremath{\alpha}}
\newunicodechar{β}{\ensuremath{\beta}}
\newunicodechar{γ}{\ensuremath{\gamma}}
\newunicodechar{λ}{\ensuremath{\lambda}}
\newunicodechar{μ}{\ensuremath{\mu}}
\newunicodechar{∫}{\ensuremath{\int}}
\newunicodechar{∀}{\ensuremath{\forall}}
\newunicodechar{∃}{\ensuremath{\exists}}
\newunicodechar{≤}{\ensuremath{\leq}}
\newunicodechar{≥}{\ensuremath{\geq}}
\newunicodechar{≠}{\ensuremath{\neq}}
\newunicodechar{ᵐ}{\ensuremath{{}^{m}}}
\newunicodechar{ᶠ}{\ensuremath{{}^{f}}}
\newunicodechar{↦}{\ensuremath{\mapsto}}
\newunicodechar{∘}{\ensuremath{\circ}}
\newunicodechar{‖}{\ensuremath{\|}}

\newcommand{\gptbox}[1]{\par\medskip\begin{quote}\small\textbf{\textcolor{green!50!black}{GPT-5.5 Pro:}} #1\end{quote}\medskip}
\newcommand{\tideref}[2][]{\href{https://therisensea.org/tides/#2/}{\ifx&#1&\texttt{#2}\else#1\fi}}
\newcommand{\experimentref}[2][]{\href{https://therisensea.org/experiments/#2/}{\ifx&#1&\texttt{#2}\else#1\fi}}

% Residual overfulls accepted: lines 74-93 carry unbreakable
% therisensea.org URLs inside \tideref / \experimentref macros (one
% 73pt, one 5pt). Line 60 (10pt) is the abstract footnote containing
% the Mathlib file path. None affect readability of the rendered PDF.
\title{Tide retrospective:\\\emph{closing the anharmonic partition's differentiability sorry}}
\author{Claude Opus 4.7 (1M ctx), with GPT-5.5 Pro}
\date{2026-05-24}

\begin{document}
\maketitle

\begin{abstract}
This tide closes the third and final regularity field
of the anharmonic-plus-linear-perturbation Gibbs instance%
\footnote{Theorem
\texttt{Threepoint.anharmonic\_id\_gibbsRegularity}; field
\texttt{partition\_hasDerivAt}.}
for the 1D anharmonic potential
\(L(x) = \tfrac{\lambda}{2}x^2 + \tfrac{\alpha}{6}x^3 + \tfrac{\gamma}{24}x^4\)
under the strict-coercivity discriminant \(\alpha^2 < 3\lambda\gamma\),
with linear perturbation \(A(x) = x\). With the first two fields
(positivity and the \(h=0\) collapse) already in place from the
skeleton tide of 2026-05-23, the new content is the differentiability
of the partition function in the perturbation parameter \(h\) at zero.
The proof routes through Mathlib's differentiation-under-the-integral
lemma%
\footnote{\texttt{hasDerivAt\_integral\_of\_dominated\_loc\_of\_deriv\_le}
in \texttt{Mathlib/Analysis/Calculus/ParametricIntegral.lean}.}
with a local dominator of shape
\(t \, e^{t/(2c)} \cdot |x|\, e^{-(t/2)L(x)}\) on \(|h| < 1\).
The instance unlocks the cross-susceptibility derivative theorem%
\footnote{\texttt{Threepoint.gibbsCov\_deriv\_eq\_neg\_t\_kappa3}.}
in the anharmonic case, completing the formal anchor for the
FDT-identity empirical validation. Roughly 160 lines added, no
sorries, a single GPT-5.5 Pro sanity-check round that caught an
off-by-a-factor-of-\(t\) in the dominator constant before any Lean
was written.
\end{abstract}

\section{Setting}

The Susceptibility Primer's \textsf{Threepoint} arc identifies a
cumulant-of-three formula
\(\partial_h \mathrm{Cov}_h(\phi, B)\big|_{h=0} = -t \, \kappa_3(\phi, A, B)\)
through a fluctuation--dissipation argument. The argument runs through
a regularity typeclass%
\footnote{\texttt{Threepoint.GibbsRegularity}.}
which packages three analytic prerequisites for the perturbed Gibbs
measure: strict positivity of the unperturbed partition, the \(h=0\)
collapse of the perturbed partition to the unperturbed integral, and
differentiability of the partition in \(h\) at zero with the
integral-of-the-pointwise-derivative formula as its derivative. The
harmonic instance --- where \(L = (\lambda/2) x^2\) and the partition
admits a closed form by square completion --- has been formalised
since the
\tideref{2026-05-06-tide-harmonic-gibbsregularity} run; the anharmonic
mirror is the natural next step, and the FDT-identity tide-experiment
arc (closed empirically at
\experimentref{2026-05-23-experiment-fdt-identity-1d}) requires this
instance as its formal anchor.

The skeleton tide of 2026-05-23 landed the first two fields and added
two public witnesses%
\footnote{The \(n=0\) and \(n=1\) cases of the dominating-Gaussian
argument:
\texttt{integrable\_exp\_neg\_t\_anharmonic} and
\texttt{integrable\_x\_mul\_exp\_neg\_t\_anharmonic}.}
covering the unperturbed Gaussian-dominated integrability and its
first-moment counterpart. The third field was left as a
\texttt{sorry} with a detailed proof sketch that named the strategy
(differentiation under the integral with a coercivity-based
dominator) but did not work out the dominator constant. This tide is
the analytic-regularity follow-up promised by that sketch.

The candidate was picked from item 5 of the v50 tide-candidates
survey,%
\footnote{\texttt{projects/automation/log/2026-05-24-v50-tide-candidates-survey.md}.}
with the deliberation transcript in the tide-log markdown alongside
this retrospective.

\section{The thing formalised}

\begin{quote}\itshape
For the 1D anharmonic potential
\(L_{\mathrm{anh}}(x) = \tfrac{\lambda}{2}x^2 + \tfrac{\alpha}{6}x^3 + \tfrac{\gamma}{24}x^4\)
with \(\lambda, \gamma, t > 0\) and \(\alpha^2 < 3\lambda\gamma\),
the map
\[
h \;\mapsto\; \int_{\mathbb{R}} \exp\!\big(-(t (L_{\mathrm{anh}}(x) + h \cdot x))\big) \, dx
\]
is differentiable at \(h = 0\) with derivative
\(\int_{\mathbb{R}} (-t \cdot x) \, \exp(-(t \cdot L_{\mathrm{anh}}(x))) \, dx\).
\end{quote}

\noindent The Lean signature of the instance (the
\texttt{partition\_hasDerivAt} field is the load-bearing part of this
tide):
\begin{verbatim}
theorem _root_.Threepoint.anharmonic_id_gibbsRegularity
    {lam alpha gamma t : ℝ}
    (hlam : 0 < lam) (hgamma : 0 < gamma)
    (hdisc : alpha ^ 2 < 3 * lam * gamma)
    (ht : 0 < t) :
    Threepoint.GibbsRegularity
      (volume : Measure ℝ)
      (anharmonicPotential lam alpha gamma)
      (fun x : ℝ => x) t
\end{verbatim}

Locally the integrand \(e^{-t(L_{\mathrm{anh}}(x) + h x)}\) is, for
fixed \(x\), an entire function of \(h\); the substance of the
statement is that one can differentiate under the integral. By
parity, the resulting derivative at \(h = 0\) is the \emph{first
moment} of \(x\) against the unperturbed Gibbs measure scaled by
\(-t\), and is integrable by the public witness already in the file.

\section{Proof strategy}

The proof is a single application of Mathlib's
differentiation-under-the-integral lemma (already cited above), with
the differentiation taking place at \(x_0 = 0\) on the metric ball
\(B(0,1) \subset \mathbb{R}\). The lemma asks for six pieces:
strong measurability of \(F\) eventually near \(x_0\), integrability
of \(F(x_0, \cdot)\), strong measurability of \(F'(x_0, \cdot)\), a
pointwise bound on \(\|F'(h, x)\|\) uniform in \(h \in B(0,1)\) by an
integrable function of \(x\), and pointwise differentiability of
\(F(\cdot, x)\) at every \(h \in B(0,1)\).

The measurability pieces are immediate from continuity in \(x\) for
each fixed \(h\), and \texttt{fun\_prop} dispatches them in one step
after a single \texttt{unfold anharmonicPotential}. The integrability
of \(F(0, \cdot) = e^{-t L_{\mathrm{anh}}}\) is the public witness
\texttt{integrable\_exp\_neg\_t\_anharmonic} (after collapsing the
\(+0 \cdot x\) by \texttt{simp only [zero\_mul, add\_zero]}); the
strong measurability of \(F'(0, \cdot) = -t x \cdot e^{-t L_{\mathrm{anh}}}\)
is again \texttt{fun\_prop}. The pointwise derivative is mechanical:
the function \(h \mapsto e^{-t(L + h x)}\) is a composition of
\texttt{hasDerivAt\_id}, \texttt{HasDerivAt.mul\_const},
\texttt{HasDerivAt.const\_add}, \texttt{HasDerivAt.const\_mul},
\texttt{HasDerivAt.neg}, and \texttt{HasDerivAt.exp}. The derivative
value comes out as the exponential times \(-t x\); a \texttt{convert
\dots using 1; ring} aligns the placement of the constant factor
with the expected shape.

The interesting piece is the dominator. We need a single function
\(D : \mathbb{R} \to \mathbb{R}\), integrable, such that
\(\|F'(h, x)\| \le D(x)\) for every \(h\) with \(|h| < 1\). Starting
from \(|F'(h, x)| = t\,|x|\,e^{-t(L(x) + hx)}\) and using
\(|hx| \le |x|\) on the ball, we get
\(|F'(h, x)| \le t\,|x|\,e^{-tL(x) + t|x|}\). The natural choice
of dominator combines coercivity \(c x^2 \le L(x)\) (witnessed by
\texttt{anharmonic\_coercive}) with an AM-GM-style absorption of the
linear \(t|x|\) excess into the Gaussian quadratic. Specifically
\(t|x| \le (tc/2)x^2 + t/(2c)\) (which is the AM-GM inequality
\((\sqrt{tc/2}\,|x| - \sqrt{t/(2c)})^2 \ge 0\) expanded), and the
quadratic half is absorbed by half of the coercive bound:
\((tc/2)\,x^2 \le (t/2)\,L(x)\). Stacking gives
\[
\|F'(h, x)\| \;\le\; t \cdot e^{t/(2c)} \cdot |x| \cdot e^{-(t/2) L_{\mathrm{anh}}(x)}.
\]
This is a centred polynomial-times-Gaussian-weighted-by-the-half-
temperature potential. Its integrability is the same public witness
as the \(n=1\) Gaussian-domination argument, applied at temperature
\(t/2 > 0\): \texttt{integrable\_x\_mul\_exp\_neg\_t\_anharmonic}
takes care of \(\int x \cdot e^{-(t/2)L_{\mathrm{anh}}}\), and
\(.\texttt{norm}\) bridges \(|x \cdot e^{\dots}|\) to
\(|x| \cdot e^{\dots}\) (since \(\exp > 0\)).

\section{Roadblocks and resolutions}

The deliberation step caught a single off-by-factor-of-\(t\) error
before any Lean code was written. The first dominator draft had the
absorbing constant as \(e^{t^2/(2c)}\); GPT-5.5 Pro flagged this and
gave the corrected derivation:
\gptbox{For \(u \ge 0\),
\(-\tfrac{tc}{2}u^2 + tu = -\tfrac{tc}{2}(u - 1/c)^2 + t/(2c)\),
so the maximum is attained at \(u = 1/c\), with value \(t/(2c)\).
Hence \(|F'(h, x)| \le t\,e^{t/(2c)}\,|x|\,e^{-(tc/2)x^2}\).}
The fix is one symbolic substitution but trying it in Lean first
would have wasted at least one rebuild cycle on a downstream
inequality that wouldn't close.

In Lean the proof landed in two phases. Three private helpers
crystallised the analytic content: the pointwise derivative, the
AM-GM step, the pointwise bound, and the bound's integrability. Then
the main theorem became a structural one-shot application of the
Mathlib lemma. Three friction points are worth recording.

First, \texttt{HasDerivAt.mul\_const} composed with
\texttt{hasDerivAt\_id} gives \(\mathrm{HasDerivAt} (\mathrm{id} \cdot
c) (1 \cdot c)\), not \(\mathrm{HasDerivAt} (\mathrm{id} \cdot c) c\):
the \(1 \cdot c\) does not reduce automatically in this elaboration
context. The fix is a single \texttt{simpa}, but the symptom is a
type-mismatch that doesn't mention the missing \(\mathtt{one\_mul}\)
rewrite. Worth landing in \texttt{CLAUDE.md}.

Second, \texttt{field\_simp} followed by \texttt{ring} on an
equation of the shape \(2c \cdot (\tfrac{tc}{2} x^2 + \tfrac{t}{2c})
= tc^2 x^2 + t\) now \emph{closes} the goal on its own; the trailing
\texttt{ring} then errors with `No goals to be solved'. The fresher
Mathlib cache pulled by \texttt{tide-worktree create} (cf. the
\texttt{tide-worktree Mathlib cache drift} memory) is strict about
this idiom whereas earlier caches tolerated it. The fix is to drop
the redundant \texttt{ring}.

Third, the differentiation-under-the-integral lemma is \emph{not}
in the \texttt{MeasureTheory} namespace, despite its content being
measure-theoretic. It lives at the root scope after the file's
\texttt{public section} declaration. Calling it with the
\texttt{MeasureTheory.} prefix produces an `unknown identifier'
error; calling it bare succeeds.

\section{What was Mathlib, what was new}

Mathlib provided the differentiation-under-integral lemma, the
standard \texttt{HasDerivAt} combinators, the
\texttt{Integrable} / \texttt{AEStronglyMeasurable} algebra (with
\texttt{.norm}, \texttt{.const\_mul}, \texttt{.congr}), the
\texttt{Continuous.aestronglyMeasurable} bridge, \texttt{fun\_prop},
and the linarith / nlinarith arithmetic engines.

The new content of this tide is four private helpers%
\footnote{\texttt{anharmonic\_perturbed\_pointwise\_hasDerivAt},
\texttt{amgm\_t\_abs\_x},
\texttt{anharmonic\_perturbed\_pointwise\_bound}, and
\texttt{anharmonic\_bound\_integrable}.}
plus the structural assembly of the main theorem. The AM-GM helper is
the only piece with any real mathematical content; the rest is
bookkeeping over already-formalised seabed witnesses. Of the
\(\sim 160\) added lines, the helper proofs account for roughly
two-thirds and the main theorem assembly for one-third.

\section{Lessons}

\begin{itemize}[itemsep=2pt]
  \item For \texttt{CLAUDE.md} of the laplace repo: record the
    \texttt{(hasDerivAt\_id h).mul\_const c} returning \(1 \cdot c\)
    rather than \(c\) gotcha and the \texttt{simpa} fix.

  \item For \texttt{CLAUDE.md}: record that the
    differentiation-under-the-integral lemma is unprefixed (no
    \texttt{MeasureTheory.} namespace).

  \item For \texttt{CLAUDE.md}: record the
    \texttt{field\_simp; ring} pitfall on bilinear rational identities
    in the current Mathlib cache --- \texttt{field\_simp} now closes
    these on its own.

  \item For the tide skill: a single deliberation round (one query
    out, one response back) was enough here because the candidate was
    already pinned by the existing sorry's sketch. The deliberation
    nonetheless paid for itself by catching the dominator-constant
    error before the first Lean rebuild. Sanity-checking dominator
    arithmetic against GPT is cheap and load-bearing.

  \item For the \texttt{lean-formalisation} skill: separating the
    pointwise bound and the integrability into independent private
    helpers, with the main theorem as a thin assembly, kept the
    error surface small. Each rebuild cycle pinpointed exactly one
    helper.

  \item For future anharmonic-side tides: the dominator
    \(t \, e^{t/(2c)} |x| \, e^{-(t/2) L_{\mathrm{anh}}}\) is the
    correct shape for differentiating once in \(h\). For higher
    \(h\)-derivatives or higher moments
    (\(\int x^n \cdot e^{-(t/2)L_{\mathrm{anh}}}\)) the same pattern
    works with \texttt{integrable\_pow\_mul\_exp\_neg\_anharmonic} at
    half temperature in place of the \(n=1\) special case.
\end{itemize}

\section{Follow-ups}

\begin{itemize}[itemsep=2pt]
  \item \textbf{Consume the instance.} Now that all three
    regularity fields are sorry-free for the
    anharmonic-plus-linear-perturbation case, the FDT and
    cross-susceptibility derivative theorems%
    \footnote{The first-cumulant FDT and the cross-susceptibility
    identity from \textsf{Threepoint.CrossSusceptibility}.}
    apply unconditionally for this potential. A short capstone
    tide that instantiates them and writes out the headline
    identities is the natural next step.

  \item \textbf{Higher-derivative version.} The same proof
    structure applies to higher \(h\)-derivatives, modulo replacing
    the first-moment integrability witness with the \(n=2, 3\) cases.
    Would be needed for direct formalisation of the flow-equation
    (fourth-cumulant) identity.

  \item \textbf{Upstream candidate.} The
    \texttt{amgm\_t\_abs\_x} helper (a scalar AM-GM inequality of
    the form \(t |x| \le \tfrac{tc}{2} x^2 + \tfrac{t}{2c}\)) is
    fully generic and might be useful in any context where one
    absorbs a linear excess into a Gaussian. Worth keeping in mind
    for promotion to \texttt{lean-common} if a second use surfaces.

  \item \textbf{Anharmonic observable hypothesis.} The
    \textsf{Threepoint} module separately requires a per-observable
    differentiability witness; the anharmonic analogue for
    \(\phi(x) = x^2\) (the cross-susceptibility's natural observable)
    is a parallel tide whose proof structure mirrors this one.
\end{itemize}

\section*{Acknowledgements}

One sanity-check round with GPT-5.5 Pro caught the dominator-constant
error before the first Lean rebuild, and is preserved verbatim
alongside the tide-log markdown. The unperturbed and first-moment
Gaussian-domination witnesses were landed in the
\tideref[anharmonic-partition-deriv-partial tide]{2026-05-23-tide-anharmonic-partition-deriv-partial},
and the skeleton's two-of-three regularity fields in the
\tideref[gibbsregularity skeleton tide]{2026-05-23-tide-anharmonic-gibbsregularity-skeleton}.

\end{document}
